\section{The exact carrier quotient for an arbitrary incomplete table}\label{sec:mainquotient}

\subsection{Row-stable congruences}

For $\theta\in\Con(A)$ let $Q=A/\theta$ and project every supplied row occurrence to $Q$.  For each named $b$, form the raw row pushout \eqref{eq:rawpo} over $Q$.

\begin{definition}[Row stability]
A congruence $\theta\in\Con(A)$ is \emph{row-stable} if, for every named $b\in B$, the distinguished carrier leg of the corresponding raw row pushout over $A/\theta$ is injective.  Equivalently, every row satisfies the two conditions of Theorem~\ref{thm:raw}.
\end{definition}

\begin{theorem}[Least row-stable carrier quotient]\label{thm:leaststable}
For the pure row fragment,
\begin{equation}\label{eq:leaststable}
\boxed{
\theta_E^A
=
\min\{\theta\in\Con(A):\theta\text{ is row-stable}\}.
}
\end{equation}
In particular, the full carrier congruence forced by an arbitrary incomplete action table is determined entirely inside the finite congruence lattice of $A$.
\end{theorem}

\begin{proof}
Put $\theta=\theta_E^A$ and $Q=Q_E$.  The named carrier quotient $Q$ embeds in the initial $\A$-sort.  For every row, Proposition~\ref{prop:rowhom} and the supplied cells realize the corresponding gluing inside that initial algebra while preserving $Q$.  By the universal property of the raw pushout, its carrier leg must therefore be injective.  Thus $\theta_E^A$ is row-stable.

Conversely let $\theta$ be row-stable and put $Q=A/\theta$.  For each $b$, Theorem~\ref{thm:raw} gives a homomorphism $h_b:S_b\to Q$ whose generated kernel $\kappa_b$ satisfies the extension condition.  Factor the row copy by $\kappa_b$ and strongly amalgamate $Q/\kappa_b$ with the distinguished $Q$ along the common copy of $S_b/\ker h_b$.  Do this successively for all finitely many rows, always over the same embedded distinguished copy of $Q$.  Lemma~\ref{lem:sap} preserves that central copy at every step.

Use the resulting amalgam as the carrier sort of a model.  Interpret the operator sort exactly as $B$.  For a named $b$ and a central element $[a]\in Q$, define $\alpha(b,[a])$ to be the corresponding point of the $b$-row image; define all remaining action values arbitrarily.  Native operations are those of the amalgam.  Carrier compatibility holds on the named central inputs because each row image is a homomorphic quotient of $Q$, and the supplied cells hold by construction.  Hence this is a model of $E$ whose kernel on the named carrier is exactly $\theta$.

Every provable carrier equality must hold in this model, so $\theta_E^A\subseteq\theta$.  Since $\theta_E^A$ itself is row-stable, it is the least such congruence.
\end{proof}

Theorem~\ref{thm:leaststable} is the structural center of the paper.  It separates two effects that are conflated by direct table completion.  The quotient $\theta_E^A$ is forced by unrestricted equational semantics; only after this quotient has been taken does each row possess a canonical forced partial homomorphism, whose exact domain is given by Theorem~\ref{thm:kernelsat}.


\subsection{All rows inside the full initial algebra}

The preceding proof also supplies the separation model needed to transfer
local pushout geometry to the full presentation. This transfer is important:
a local upper bound on forcing would not by itself rule out extra
identifications caused by other rows.

\begin{theorem}[Exact global row geometry]\label{thm:globalrows}
Let $Q=Q_E$ and let $\iota:Q\hookrightarrow\mathbb I_E^{\A}$ be the named
carrier embedding. For each nonempty row, form $S_b,h_b,K_b,\kappa_b$ over
$Q$. Then
\begin{align}
\ker\bar j_b&=\kappa_b,\label{eq:globalrowkernel}\\
\bar j_b(Q)\cap\iota(Q)&=\iota(h_b(S_b)),\label{eq:globalrowintersection}\\
D_b&=S_b^{\kappa_b}.\label{eq:globaldomain}
\end{align}
If $x\equiv_{\kappa_b}s$ with $s\in S_b$, its forced value is $h_b(s)$.
The external parts of distinct named rows are disjoint in the initial
algebra. For an empty row, the row kernel is diagonal, its image is disjoint
from the carrier and all other row images, and $D_b=\varnothing$.
\end{theorem}

\begin{proof}
Theorem~\ref{thm:leaststable} makes every row protected over $Q$.
The supplied equations imply $\bar j_b(s)=\iota(h_b(s))$ on $S_b$.
Consequently $K_b\subseteq\ker\bar j_b$, and the row-homomorphism principle
implies $\kappa_b\subseteq\ker\bar j_b$. This proves all equalities claimed
by the theorem, including the values on $S_b^{\kappa_b}$.

For exactness, construct a model as in the second half of the proof of
Theorem~\ref{thm:leaststable}. At each step attach $Q/\kappa_b$ by a strong
amalgam along the embedded $S_b/K_b$ in the central $Q$. The new row meets
the entire previously constructed algebra only along that subalgebra of
$Q$. Thus its kernel is exactly $\kappa_b$, its carrier intersection is
exactly $h_b(S_b)$, and it adds no intersection between external row parts.
Attach an empty row by a disjoint coproduct. Interpreting the action on
central inputs by these row maps gives a model of all of $E$.
Every pair excluded by the theorem is separated in this one model, so it
cannot become equal in the initial algebra.
\end{proof}

\subsection{A finite feedback computation}

The theorem also gives a canonical closure procedure.  For $\theta\in\Con(A)$ let $Q_\theta=A/\theta$, form every raw row pushout over $Q_\theta$, and write
\[
\chi_{b,\theta}=\ker(m_{b,\theta})\in\Con(Q_\theta).
\]
Let $\pi_\theta:A\twoheadrightarrow Q_\theta$ and define
\begin{equation}\label{eq:Phi}
\Phi(\theta)
=
\bigvee_{b\in B}\pi_\theta^{-1}(\chi_{b,\theta}).
\end{equation}
For an unspecified row, $\chi_{b,\theta}=\Delta_{Q_\theta}$.

\begin{proposition}[Pushout feedback operator]\label{prop:Phi}
The operator $\Phi:\Con(A)\to\Con(A)$ is inflationary and monotone.  Its fixed points are exactly the row-stable congruences.  Consequently
\begin{equation}\label{eq:lfp}
\boxed{\theta_E^A=\operatorname{lfp}(\Phi).}
\end{equation}
Starting with $\theta^{[0]}=\Delta_A$ and $\theta^{[r+1]}=\Phi(\theta^{[r]})$, at most $|A|-1$ strict steps occur.
\end{proposition}

\begin{proof}
Each pullback $\pi_\theta^{-1}(\chi_{b,\theta})$ contains $\theta$, giving inflationarity.  If $\theta\subseteq\varphi$, the quotient map $A/\theta\twoheadrightarrow A/\varphi$ induces a morphism between the two raw row spans and hence between their pushouts.  Equality created in the carrier leg at level $\theta$ therefore remains equality after passage to $\varphi$, which gives monotonicity after pulling back to $A$.

The join in \eqref{eq:Phi} equals $\theta$ precisely when every $\chi_{b,\theta}$ is diagonal, i.e. precisely when every raw row pushout protects the carrier.  The fixed-point statement now follows from Theorem~\ref{thm:leaststable}.  In a strict step the number of congruence classes decreases by at least one, so a finite $A$ allows at most $|A|-1$ strict steps.
\end{proof}

\begin{proposition}[Sharpness of the feedback bound]\label{prop:sharpfeedback}
For every $n\ge2$ there is an $n$-element unary carrier and an incomplete action presentation for which the iteration of Proposition~\ref{prop:Phi} has exactly $n-1$ strict steps.
\end{proposition}

\begin{proof}
Let $A_n=\{0,1,\ldots,n-1\}$ with one unary operation $f(x)=0$ for all $x$.  Take distinct operators $b_1,\ldots,b_{n-1}$ and supply exactly
\[
\alpha(b_i,i-1)=i
\qquad(1\le i<n).
\]
At the first stage the $b_1$ row asks for $h(0)=1$.  A protected homomorphic row would satisfy
\[
1=h(0)=h(f(0))=f(h(0))=f(1)=0,
\]
so $0\equiv1$ is forced.  After quotienting by this equality, the source of the $b_2$ row is the zero class and the same argument forces $[0]=[2]$.  To check that no other collapse occurs in the same round, suppose the
current blocks are $\{0,\ldots,k\}$ and the remaining singletons. Rows with
$i\le k$ already map the zero class to itself. Row $i=k+1$ forces exactly
its target into the zero class. For $i>k+1$, the supplied input generates
the two-element subalgebra consisting of the zero class and $[i-1]$;
the prescribed extension sends these to the zero class and $[i]$.
It is an injective homomorphism and hence protected by
Theorem~\ref{thm:raw}. Because the update is simultaneous, these later rows
contribute no additional equality in that round. Induction therefore gives
exactly the block $\{0,\ldots,r\}$ after $r$ strict rounds, and total
collapse first occurs after $n-1$ rounds.
\end{proof}

\subsection{Two endpoint corollaries}

The earlier special cases now become immediate consequences.

\begin{corollary}[Annihilator row]\label{cor:ann}
Assume designated elements $0_A,0_B$ are present and that $\Gamma$ consists only of the complete row
\[
\alpha(0_B,a)=0_A
\qquad(a\in A).
\]
All other rows are unspecified. Then the exact carrier congruence is
\begin{equation}\label{eq:thetaann}
\theta_{\rm Ann}
=
\Cg_A\{(0,f^A(0,\ldots,0)):f\in\Sigma_A\}.
\end{equation}
\end{corollary}

\begin{proof}
On a quotient $Q$, the annihilator row is the constant-zero map $Q\to Q$.  Since the row is complete, stability is exactly the assertion that this map is an endomorphism, i.e. that $[0]$ is closed under every positive-arity carrier operation.  The least quotient with that property is \eqref{eq:thetaann}.
\end{proof}

\begin{corollary}[Complete action table]\label{cor:complete}
Suppose one value
\[
\alpha(b,a)=\tau_b(a)
\]
is supplied for every $(b,a)\in B\times A$.  Then $\theta_E^A$ is the least native carrier congruence $\theta$ for which every $\tau_b$ descends to an endomorphism of $A/\theta$.  Equivalently, it is generated, in the expanded algebra $(A,\Sigma_A,\{\tau_b\}_{b\in B})$, by the compatibility defects
\[
\bigl(\tau_b(f(\bar a)),\ f(\tau_b(a_1),\ldots,\tau_b(a_k))\bigr).
\]
\end{corollary}

Thus the annihilator and complete-table formulas are not isolated regimes.  They are instances of the case $S_b=Q$ of the same arbitrary-partial-row theorem.

